\section{量子化：预量子化、极化与全纯截面}

\subsection{为什么量子化不平凡}

我们现在有了一个有限维、紧致、Kähler 的相空间 $\cM_J$，要把它量子化。
如果 $\cM_J$ 是 $\RR^{2n}$，量子化就是“把 $p$ 换成 $-i\hbar\partial_x$”，没有歧义。
但 $\cM_J$ 是紧致的，没有整体的“位置与动量分离”：任何试图把坐标分成 $q$ 与 $p$ 的做法
都只能局部进行。更具体地，量子态应当是相空间上的“函数”，但“函数”是什么空间？
如果取所有 $L^2$ 函数，维数是无限的，而物理上我们希望有限维（因为相空间体积有限）。

几何量子化给出了一套系统的做法，分三步：\textbf{预量子化}、\textbf{极化}、\textbf{半形式修正}。
这三步分别解决“线丛从哪来”“态是什么”“能级怎么移”三个问题。

\subsection{第一步：预量子化线丛}

相空间 $(\cM,\omega)$ 的预量子化，是找一个复线丛 $L\to\cM$，带一个酉联络 $\nabla$，
使得它的曲率是
\[
F_\nabla=-\ii\,\omega
\]
（这里 $\omega$ 按“每单位体积一个态”的归一化，见下）。这样的线丛存在的充要条件是
\begin{equation}
\frac{[\omega]}{2\pi}\in H^2(\cM;\ZZ). \tag{4.1}
\end{equation}
这就是“整性条件”。满足它之后，$L$ 上的截面就构成预量子 Hilbert 空间。

在我们的情形里，$\cM$ 上的辛形式由 (3.2) 给出：
\[
\omega=\frac{k}{4\pi}\int_\Sigma\Tr\left(\delta A\wedge\delta A\right).
\]
注意 $\omega$ 本身正比于 $k$。整性条件 (4.1) 于是恰好要求 $k$ 取整数——这与第二讲从大规范变换
得到的 $k\in\ZZ$ 是同一件事的两种说法。这一点很重要：\textbf{量子化的自洽性自动包含了耦合常数的量子化}。

那么线丛 $L$ 具体是什么？答案是行列式线丛。$\cM$ 是平坦联络的模空间，而平坦联络可以看成一族
椭圆算子（$\bar\partial_A$）的参数的模空间；这族算子的行列式构成 $\cM$ 上的一个线丛，
Quillen 给了它一个自然的度量与联络，其曲率正是上面的 $\omega$。这个构造的好处是：
\textbf{线丛是先验存在的}，所以整性条件自动满足，不需要额外验证。

\subsection{第二步：极化}

预量子 Hilbert 空间是 $L$ 的 $L^2$ 截面空间。它太大了：$\cM$ 的维数是 $2n$ 时，
这个空间是“$n$ 个坐标的函数”，而量子态应该只依赖一半的变量。
物理上，这对应“不能同时指定位置与动量”。

解决办法是选一个\textbf{极化}：复化切丛 $T\cM\otimes\CC$ 的一个子丛 $P$，满足
\begin{enumerate}[leftmargin=*]
  \item $P$ 是拉格朗日的：$\omega|_P=0$；
  \item $P$ 是可积的：对合（即 $P$ 的截面在 Lie 括号下封闭）；
  \item $\dim_\CC P=n$。
\end{enumerate}
量子态定义为“沿 $P$ 协变常数的截面”：对 $P$ 的任意截面 $X$，要求 $\nabla_Xs=0$。
这样，态只依赖 $P$ 的叶空间上的变量，维数就从无限降下来。

最常见的两种极化是：
\begin{itemize}[leftmargin=*]
  \item \textbf{实极化}：$P$ 是实切丛的复化，对应“位置空间”或“动量空间”。在 $\RR^{2n}$ 上取
        $P=\mathrm{span}\{\partial/\partial p_i\}$，得到的态只依赖 $x$。
  \item \textbf{Kähler 极化}：$P=T^{0,1}\cM$。此时“沿 $P$ 协变常数”就是“全纯”，
        量子态是全纯截面。
\end{itemize}
我们的 $\cM_J$ 是 Kähler 的，自然取 Kähler 极化。于是
\begin{equation}
\cH_\Sigma=H^0(\cM_J,\,L^{\otimes k}), \tag{4.2}
\end{equation}
即线丛 $L$ 的 $k$ 次幂的全纯截面空间。这里出现 $L^{\otimes k}$ 而不是 $L$，
是因为 $\omega$ 正比于 $k$：整性条件要求线丛的曲率是 $\omega$，而 $\omega=k\omega_0$，
所以线丛取 $L^{\otimes k}$，其中 $L$ 的曲率是 $\omega_0$。

\subsection{维数与半经典极限}

(4.2) 的维数由 Riemann--Roch（或指标定理）给出。对 $SU(N)$、亏格 $g$ 的情形，
$k\to\infty$ 时
\[
\dim H^0(\cM_J,L^{\otimes k})\sim \frac{\mathrm{Vol}(\cM)}{(2\pi)^n n!}\,k^n,
\]
即“每单位体积一个态”的 Weyl 律，其中 $2n=\dim\cM$。

最简单地，取 $g=1$、$G=SU(2)$：此时 $\cM_J$ 是枕头套 $T^2/\{\pm1\}$，
而 $\dim H^0(\cM_J,L^{\otimes k})=k+1$——这正是环面上的 theta 函数空间维数。
$k=1$ 时给出 $2$，$k=2$ 时给出 $3$，与第八讲算 $S^3$ 配分函数时用到的数一致。

值得把 $g=1$ 的情形算清楚，因为它包含了几何量子化最朴素的所有要素。取环面
$\Sigma=\CC/(\ZZ+\tau\ZZ)$。度为 $k$ 的全纯线丛的标准截面由 theta 函数给出：
\[
\vartheta_{j}(z)=\sum_{m\in\ZZ}\exp\left(\pi\ii k(m+\tfrac jk)^2\tau
+2\pi\ii k(m+\tfrac jk)z\right),\qquad j=0,1,\dots,k-1,
\]
共 $k$ 个。但正确的量子化还要乘上半形式因子 $K^{1/2}$；对 $SU(2)$ 它是 $L^{-2}$，
于是实际空间是 $H^0(L^{k}\otimes K^{1/2})$，维数为 $k+1$。这正好等于
“有效 level $k+2$ 的可积表示个数减一”，与 (4.3) 的平移一致。
这个例子说明：半形式不是一个可选的修饰，它改变维数。

\subsection{第三步：与复结构无关吗？}

一个自然的问题是：$\cH_\Sigma=H^0(\cM_J,L^{\otimes k})$ 依赖复结构 $J$，
但物理上 Hilbert 空间不应该依赖“怎么切”。这两个说法如何统一？

答案是：当 $J$ 变化时，$H^0(\cM_J,L^{\otimes k})$ 拼成 Teichmüller 空间上的一个向量丛，
而这个丛带有\textbf{射影平坦}联络（Hitchin 联络）。所谓射影平坦，是指平行移动只在一个
整体相位上有歧义；沿闭路平行移动给出一个射影变换，其单值性正是映射类群在 $\cH_\Sigma$ 上的作用。
因此严格地说，$\cH_\Sigma$ 不是一个“与 $J$ 无关的向量空间”，而是一个“射影平坦丛的纤维”；
只有射影化之后的 $\PP\cH_\Sigma$ 才是典范的。

这一点在第五讲之后会变得关键：共形块正是这个射影平坦丛的纤维，
而映射类群的作用给出辫群表示，从而给出纽结不变量。

\subsection{半形式修正与能级平移}

还有第三步：半形式修正。Kähler 极化要求我们取 $L^{\otimes k}$ 的截面，
但严格地，量子化还需要一个“半形式”因子 $K^{1/2}$（$K$ 是 $\cM_J$ 的典范丛）。
对 $\cM_J$ 而言，$K$ 与 $L$ 的某个幂有关：$K=\mathrm{Det}^{-2h^\vee}$，
因此 $K^{1/2}=L^{-h^\vee}$。把这个因子乘进去，等价于把 $k$ 换成
\begin{equation}
k\;\longrightarrow\;k+h^\vee. \tag{4.3}
\end{equation}
对 $SU(N)$，$h^\vee=N$，所以 $k\to k+N$；对 $SU(2)$，$k\to k+2$。

(4.3) 是一个“能级平移”。它在物理上来自标架反常，在数学上来自半形式。
Witten 的论文在正文里只用 $L^{\otimes k}$，把这个修正略去了；但所有需要
$q=\exp(2\pi\ii/(k+N))$ 的公式，都已经隐含了它。读者在比较不同文献时，
必须注意对方的 $k$ 是“裸”的还是“已平移”的。

\subsection{Hitchin 联络}

“射影平坦”这句话可以写成一个具体的微分方程。设 $t$ 是 Teichmüller 空间的复坐标（复结构参数），
$s_t\in H^0(\cM_J,L^{\otimes k}\otimes K^{1/2})$ 是一族截面。Hitchin 联络是
\begin{equation}
\nabla_{\partial_t}s=\partial_t s+\frac{1}{k+h^\vee}\left(\Theta_t\right)s, \tag{4.4}
\end{equation}
其中 $\Theta_t$ 是由 Beltrami 微分与模空间上的二阶基本形式构造的算符（具体形式依赖符号约定）。
这个联络的曲率是纯中心的：$[\nabla_{\partial_t},\nabla_{\partial_{\bar t}}]=\frac{\ii}{2}\omega_{\mathrm{WP}}$，
其中 $\omega_{\mathrm{WP}}$ 是 Weil--Petersson 形式。因此平行移动只差一个整体相位，
这正是“射影平坦”。

若把 $t$ 取成带标记点曲面的模参数，并把 $K^{1/2}$ 的贡献显式写出，(4.4) 就变成 KZ 方程 (6.5)。
这是第三、四讲的三维语言与第六讲的二维语言之间最直接的一座桥：Hitchin 联络的平坦性
（三维侧要求）等价于共形块的 Ward 恒等式（二维侧要求）。

\subsection{小结}

几何量子化的三步，在我们的情形里都有具体的答案：
\begin{enumerate}[leftmargin=*]
  \item 预量子化：行列式线丛，整性自动；
  \item 极化：Kähler 极化，量子态 $=H^0(\cM_J,L^{\otimes k})$；
  \item 半形式：把 $k$ 平移成 $k+h^\vee$。
\end{enumerate}
至此，三维理论的 Hilbert 空间有了一个完全具体的描述。下一讲把纽结放进来。

\subsection*{练习}
\begin{enumerate}[leftmargin=*]
  \item 验证 (4.1) 与 $k\in\ZZ$ 的等价性（提示：$\omega$ 的周期积分）。
  \item 对 $\RR^2$ 取 $P=\mathrm{span}\{\partial/\partial p\}$，写出满足 $\nabla_Xs=0$ 的截面，并说明它只依赖 $x$。
  \item 用 Riemann--Roch 计算 $g=1$、$SU(2)$、$L^{\otimes k}$ 的截面数，并说明半形式修正如何把它变成 $k+1$。
  \item 解释 (4.4) 中的分母为什么是 $k+h^\vee$ 而不是 $k$。
\end{enumerate}

\begin{quote}
\small
\textbf{本讲要点。} 量子态是 $\cM_J$ 上 $L^{\otimes k}$ 的全纯截面；$J$ 依赖只到射影平坦；
半形式修正把 $k$ 平移为 $k+h^\vee$。
\end{quote}
